“Stttt, wees stil! De buren!”
Moeder moest weer gluren

naar de pachtboerderij.
Die lag het meest vlakbij.

“Ze bespieden ons steeds.
Ze onthulden ons reeds

veel te goed door jullie
klikkende spraakherrie.”

“De buren zijn er niet”
zei de kleinste deugniet.

“Jawel! Elke ochtend
verschuift eerst die rotvent

zijn jute gordijnen,
voordat ze verdwijnen.”

Moeder wist het zeker
en drukte een beker

gewend en behoedzaam,
strak tegen een loodraam.

Emmy keek moeder aan.
Die geloofde haar waan.

“Gisternacht was het raak.
Toen hoorde ik gekraak

rondom de boerderij.
Dan verduisteren zij

wat van onze voorraad.
Zij zijn het pure kwaad.

Geloof me maar Emmy!”
Tijdens een nachtmerrie

had moeder ‘s nachts geschreeuwd.
Het had voorwaar gesneeuwd,

maar men kon niets horen
en vond geen voetsporen

wat haar relaas bewees.
Toch volharde haar vrees

dat ze werd bestolen
van haar graan en kolen,

hoewel pa teruggaf
dat hij er wat vanaf

had gehaald ter gebruik.
Moeder bleef bij misbruik

en haar paranoia
verhulde een trauma,

die ze soms verhaalde,
maar dan vlug afdwaalde

en naar de buren bracht.
Ze had niet de geestkracht

haar pijn te omarmen
en die te verarmen.

De waan werd diverser.
De nijd werd perverser.
